\documentclass[10pt]{article}
\usepackage[a4paper,margin=19mm,top=16mm,bottom=16mm]{geometry}
\usepackage{booktabs}
\usepackage{array}
\usepackage{tabularx}
\usepackage{xcolor}
\usepackage{titlesec}
\usepackage{fancyhdr}
\usepackage[hidelinks]{hyperref}
\usepackage[T1]{fontenc}
\usepackage{lmodern}
\definecolor{desk}{HTML}{1F3A5F}
\definecolor{rule}{HTML}{9FB3C8}
\pagestyle{fancy}
\fancyhf{}
\renewcommand{\headrulewidth}{0.4pt}
\fancyhead[L]{\footnotesize\color{desk}PCD --- Payments Coverage Desk}
\fancyhead[R]{\footnotesize\color{desk}\thepage}
\fancyfoot[C]{\scriptsize\color{black!55}Generated by the PCD controller; every figure is bound to a metric\_id and an SEC accession.}
\titleformat{\section}{\normalfont\bfseries\color{desk}\large}{}{0pt}{}
\titlespacing*{\section}{0pt}{7pt}{3pt}
\setlength{\parindent}{0pt}
\setlength{\parskip}{3.2pt}
\renewcommand{\arraystretch}{1.12}
\newcommand{\smallnote}[1]{\footnotesize\color{black!65}#1}

\begin{document}
\begin{center}
{\color{desk}\rule{\textwidth}{0.9pt}}\\[3pt]
{\LARGE\bfseries\color{desk} Payments coverage: FY2025 annual update}\\[2pt]
{\small Payments coverage (sell-side) \quad|\quad V, MA, PYPL, FI, GPN}\\[1pt]
{\footnotesize Information cutoff 2026-04-30 \quad|\quad Analyst of record: Christian (analyst of record)}
\\[1pt]{\color{desk}\rule{\textwidth}{0.4pt}}
\end{center}
\vspace{-4pt}
{\small\itshape Mastercard grew fastest and Visa turned the most revenue into cash; the two networks again out-earned the processors on every margin line.}
\vspace{2pt}
\section*{What we compared}
We standardised the annual filings of five listed payments businesses from SEC EDGAR XBRL, using each issuer's latest annual report filed on or before the information cutoff. Revenue, operating income, net income, operating cash flow and capital expenditure are taken as reported in the filings; every ratio below is computed from unrounded figures rather than from the rounded numbers printed here. The fiscal year-ends do not align.\textsuperscript{\dag}\par
\begin{center}\footnotesize
\begin{tabular}{lrrrrr}
\toprule
\textbf{} & \textbf{Revenue (\$bn)} & \textbf{Growth} & \textbf{Operating margin} & \textbf{Net margin} & \textbf{FCF conversion} \\
\midrule
Visa (Sep-25) & 40.0 & 11.3\% & 60.0\% & 50.1\% & 107.6\% \\
Mastercard (Dec-25) & 32.8 & 16.4\% & 57.6\% & 45.6\% & 114.6\% \\
PayPal (Dec-25) & 33.2 & 4.3\% & 18.3\% & 15.8\% & 106.3\% \\
Fiserv (Dec-25) & 21.2 & 3.6\% & 27.5\% & 16.4\% & 123.5\% \\
Global Payments (Dec-25) & 7.7 & -0.4\% & 22.8\% & 18.2\% & 145.6\% \\
\bottomrule
\end{tabular}
\\[1pt]{\scriptsize\color{black!65}Latest reported fiscal year. Revenue in \$bn; growth is year on year; margins and conversion are ratios of reported values.}
\end{center}
\section*{Growth}
Mastercard grew revenue 16.4\% to \$32.8bn, the fastest in the group, and its three-year compound rate of 13.8\% is the only one above ten percent. Visa added 11.3\% to \$40.0bn, a three-year compound rate of 10.9\%. PayPal and Fiserv grew 4.3\% and 3.6\% respectively, with compound rates of 6.4\% and 6.1\%. Global Payments reported -0.4\% at \$7.7bn; that figure is on a post-divestiture basis, so the two-year compound rate of 2.2\% is the only growth measure comparable with the business as it now stands.\par
\section*{Profitability and cash generation}
The networks remain a structurally different business from the processors. Visa earned an operating margin of 60.0\% and a net margin of 50.1\%, Mastercard 57.6\% and 45.6\%. Global Payments (22.8\%) and Fiserv (27.5\%) sit in the twenties and PayPal is lowest at 18.3\%. Cash conversion is uniformly strong: every company converted more than its reported net income into free cash flow in the year, from 106.3\% at PayPal to 145.6\% at Global Payments, with Visa generating \$21.6bn and Mastercard \$17.2bn of free cash flow.\par
\section*{Valuation}
Our equity DCF discounts three years of free cash flow at a cost of equity between 8.9\% and 11.0\%, and assumes terminal growth of 2.5\%. On that basis Visa is worth \$391.0bn of equity and Mastercard \$315.2bn; per share, Mastercard \$347.92, Fiserv \$127.42, Global Payments \$119.91 and PayPal \$83.79. Visa's per-share value is withheld: the company does not tag a diluted share count in its XBRL filings, and we do not estimate an input we cannot source.\par
\begin{center}\footnotesize
\begin{tabular}{lrrrrr}
\toprule
\textbf{} & \textbf{Equity value (\$bn)} & \textbf{Per share (\$)} & \textbf{Cost of equity} & \textbf{Terminal growth} & \textbf{Terminal share} \\
\midrule
Visa & 391.0 & n/a & 8.9\% & 2.5\% & 83.9\% \\
Mastercard & 315.2 & 347.92 & 9.4\% & 2.5\% & 83.2\% \\
PayPal & 81.1 & 83.79 & 11.0\% & 2.5\% & 79.5\% \\
Fiserv & 70.0 & 127.42 & 9.4\% & 2.5\% & 82.5\% \\
Global Payments & 29.0 & 119.91 & 9.4\% & 2.5\% & 81.1\% \\
\bottomrule
\end{tabular}
\\[1pt]{\scriptsize\color{black!65}Equity DCF at the assumption file dated with this case. Per-share values use reported diluted weighted-average shares; the terminal share is the fraction of value coming from the terminal assumption.}
\end{center}
Between 79.5\% and 83.9\% of each valuation comes from the terminal assumption. These are therefore consistency checks on the assumptions rather than price targets, and no market price is used anywhere in this note.\par
\vspace{1pt}{\footnotesize\color{black!70}\dag\ Visa's fiscal year ends in September; the other four end in December. The comparison is between each company's latest annual report, not a common calendar year, and Visa's figures therefore cover a period ending roughly nine months earlier than the December year-ends.\\\dag\ Fiserv changed the XBRL tag behind depreciation and amortisation between FY2023 and FY2024, from DepreciationDepletionAndAmortization to OtherDepreciationAndAmortization. The reported series falls from roughly \$3.2bn to \$1.7bn with no economic event behind it, so EBITDA-based measures are not used anywhere in this note.\\\dag\ Global Payments re-presented FY2023 and FY2024 after its divestitures: FY2023 revenue is \$9.65bn as first filed and \$7.38bn as restated in the FY2025 annual report. Growth rates spanning the recast are not like-for-like, so only the FY2023-FY2025 window is used.\\\dag\ Operating income for FY2023 and FY2024 was re-presented on the same post-divestiture basis as revenue; the same window restriction applies to any margin comparison.\\\dag\ Cash flow from operations for FY2023 and FY2024 was re-presented following the divestitures; the FY2022 value remains on the pre-divestiture basis.\\\dag\ Depreciation and amortisation for FY2023 and FY2024 was re-presented following the divestitures.\\\dag\ Tax expense for FY2023 and FY2024 was re-presented following the divestitures.\\\dag\ Visa does not tag a diluted share count or diluted EPS in XBRL at or before the cutoff. Its per-share intrinsic value is withheld rather than estimated; an attested cover-page share count would be required to publish one.}\par
\section*{Sources and tie-out}
{\scriptsize Every figure above is generated from the frozen snapshot of Information cutoff 2026-04-30 and is bound to a metric identifier that resolves to one of the filings below. The controller recomputed each ratio from unrounded inputs (gate G09) and re-derived the valuation from the assumption file (gate G10).}\par
{\scriptsize\begin{tabular}{@{}lll@{}}\toprule Ticker & Form & Accession (filed) \\\midrule
FI & 10-K & 0000798354-25-000047 (2025-02-20) \\
FI & 10-K & 0000798354-26-000009 (2026-02-19) \\
GPN & 10-K & 0001123360-26-000008 (2026-02-20) \\
MA & 10-K & 0001141391-25-000011 (2025-02-12) \\
MA & 10-K & 0001141391-26-000013 (2026-02-11) \\
PYPL & 10-K & 0001633917-25-000019 (2025-02-04) \\
PYPL & 10-K & 0001633917-26-000024 (2026-02-03) \\
V & 10-K & 0001403161-24-000058 (2024-11-13) \\
V & 10-K & 0001403161-25-000089 (2025-11-06) \\
\bottomrule\end{tabular}}\par
\vspace{2pt}{\scriptsize\color{black!65}AI use and approval. The draft was produced by the PCD generator role from the verified metric registry; the figures were recomputed independently by the critic role; publication requires a recorded approval from the analyst of record. Internal research exercise, not investment advice.}\par
\end{document}
